
\documentclass{article}
\usepackage{colortbl}
\usepackage{makecell}
\usepackage{multirow}
\usepackage{supertabular}

\begin{document}

\newcounter{utterance}

\centering \large Interaction Transcript for game `cladder', experiment `full\_v1.5\_default', episode 8727 with refineonly{-}Qwen3.5{-}27B.
\vspace{24pt}

{ \footnotesize  \setcounter{utterance}{1}
\setlength{\tabcolsep}{0pt}
\begin{supertabular}{c@{$\;$}|p{.15\linewidth}@{}p{.15\linewidth}p{.15\linewidth}p{.15\linewidth}p{.15\linewidth}p{.15\linewidth}}
    \# & \multicolumn{2}{c}{Player} && \multicolumn{2}{c}{Game Master} \\
    \hline

    \theutterance \stepcounter{utterance}  
    & & & \multicolumn{4}{p{0.6\linewidth}}{
        \cellcolor[rgb]{0.9,0.9,0.9}{
            \makecell[{{p{\linewidth}}}]{
                \texttt{\tiny{[P1$\langle$GM]}}
                \texttt{You are an expert in causal inference. The following question is not a typical commonsense query, but rather a meticulously designed question created by a professor specializing in causal inference, intended to assess the students' mastery of the course content.} \\
\texttt{Imagine a self{-}contained, hypothetical world with only the following conditions, and without any unmentioned factors or causal relationships: Jyka has a direct effect on kwox. Yupt has a direct effect on kwox. For those who are not jyka, the probability of kwox is 32\%. For those who are jyka, the probability of kwox is 38\%. Will jyka increase the chance of kwox?} \\
\texttt{Start your answer with 'Yes' or 'No', followed by additional reasoning or evidence to support your explanation.} \\
            }
        }
    }
     \\ \\

    \theutterance \stepcounter{utterance}  
    & \multicolumn{4}{p{0.6\linewidth}}{
        \cellcolor[rgb]{0.561,0.737,0.561}{
            \makecell[{{p{\linewidth}}}]{
                \texttt{\tiny{[P1$\rangle$GM]}}
                \texttt{Yes} \\
\\ 
\texttt{The reasoning relies on the fundamental definition of a causal effect in a setting without confounding. The problem explicitly states that \$jyka\$ has a direct effect on \$kwox\$ and that there are no unmentioned factors or causal relationships (implying no confounders). Under these conditions, the observed association between the cause and the outcome is equivalent to the causal effect.} \\
\\ 
\texttt{We are given the conditional probabilities:} \\
\texttt{1.  \$P(kwox \textbackslash{}mid \textbackslash{}text\{not \} jyka) = 32\textbackslash{}\%\$} \\
\texttt{2.  \$P(kwox \textbackslash{}mid jyka) = 38\textbackslash{}\%\$} \\
\\ 
\texttt{The causal effect of \$jyka\$ on \$kwox\$ is defined as the difference in the probability of the outcome when the cause is present versus when it is absent (the Average Treatment Effect). Calculating this difference:} \\
\texttt{\$\$ \textbackslash{}text\{Effect\} = P(kwox \textbackslash{}mid jyka) {-} P(kwox \textbackslash{}mid \textbackslash{}text\{not \} jyka) = 38\textbackslash{}\% {-} 32\textbackslash{}\% = 6\textbackslash{}\% \$\$} \\
\\ 
\texttt{Since the probability of \$kwox\$ is strictly higher when \$jyka\$ is present compared to when it is absent, and the setup guarantees that this difference is causal, \$jyka\$ increases the chance of \$kwox\$.} \\
            }
        }
    }
    & & \\ \\

    \theutterance \stepcounter{utterance}  
    & & & \multicolumn{2}{p{0.3\linewidth}}{
        \cellcolor[rgb]{0.9,0.9,0.9}{
            \makecell[{{p{\linewidth}}}]{
                \texttt{\tiny{[GM$|$GM]}}
                \texttt{Yes} \\
            }
        }
    }
    & & \\ \\

    \theutterance \stepcounter{utterance}  
    & & & \multicolumn{2}{p{0.3\linewidth}}{
        \cellcolor[rgb]{0.9,0.9,0.9}{
            \makecell[{{p{\linewidth}}}]{
                \texttt{\tiny{[GM$|$GM]}}
                \texttt{yes} \\
            }
        }
    }
    & & \\ \\

    \theutterance \stepcounter{utterance}  
    & & & \multicolumn{2}{p{0.3\linewidth}}{
        \cellcolor[rgb]{0.9,0.9,0.9}{
            \makecell[{{p{\linewidth}}}]{
                \texttt{\tiny{[GM$|$GM]}}
                \texttt{game\_result = WIN} \\
            }
        }
    }
    & & \\ \\

\end{supertabular}
}

\end{document}
